\section{Round-seventeen positive closure: a trace-free raw vector--return--roof local theorem}
\label{sec:r17-a2}

The joint law is recovered from matrix coefficients of one induced transfer
operator.  No operator trace, flat trace, Fredholm determinant, or
trace-class compensation is used.  The raw roof-frequency tail is obtained by
branchwise coarea and repeated integration by parts on genuine collision
coordinates.

\subsection{The induced operator and its fixed bundle}

Let \(R\in[0.45,0.47]\).  Choose the standard finite family of Young magnets
for the finite-horizon triangular Sinai table and let
\(F_R:Y_R\to Y_R\) be the first return to their union.  A branch \(h\) has
homology \(\kappa_h\in\mathbb Z^2\), collision return count
\(r_h\in\mathbb N\), and roof \(\tau_{R,h}\).  On a branch,
\[
 \mathcal L_{R,u,s,b}f
 =\mathcal L_R\!\bigl(
 e^{i u\cdot\kappa+i s r+i b\tau_R}f\bigr),
 \qquad (u,s)\in\mathbb T^2\times\mathbb T,
 \quad b\in\mathbb R.
\]
For the invariant density \(h_R\) and a bounded central-window functional
\(\ell_R\), define
\[
 \Phi_{n,R}^{\ell}(u,s,b)
 =\ell_R(\mathcal L_{R,u,s,b}^{n}h_R).
\]
This is exactly the Fourier transform of the inserted joint law.  All Fourier
inversion below concerns \(\Phi_{n,R}^{\ell}\), never an operator trace.

Use stable curves cut by parent singularity and homogeneity boundaries.  In a
parameter chart transport a stable curve by its parent label until a fold is
met; at the fold retain the two one-sided trace components.  Let
\(\mathcal B_R\hookrightarrow\mathcal B_{R,w}\) be the completion for the
norm
\[
 \|f\|_{\mathcal B_R}
 =\|f\|_{s,R}+b_0\|f\|_{u,R}
   +b_1\sum_{\text{parent folds}}
       \|\operatorname{Tr}^{+}f-
          \operatorname{Tr}^{-}f\|_{-1/2}.
\]
The trace term is a boundary distribution norm, not a nuclear trace.

\begin{theorem}[Uniform parent-fold Banach bundle]
\label{thm:r17-a2-bundle}
There is a finite cover of the radius interval, fixed reference spaces in each
chart, and transport maps with uniformly bounded inverses such that:
\begin{enumerate}[label=(\roman*)]
\item the branchwise Jacobians, roof multipliers, and their first four
      parameter derivatives are bounded on the transported strong/weak pair;
\item
\[
 \|\mathcal L_R^nf\|_{\mathcal B_R}
 \le C\vartheta^n\|f\|_{\mathcal B_R}
      +C\|f\|_{\mathcal B_{R,w}},
 \qquad \vartheta<1;
\]
\item branch complexity satisfies, for some \(\eta>0\),
\[
 \sum_h e^{\eta r_h}
 \bigl(\|J_h\|_{C^2}+\|\tau_{R,h}\|_{C^4}
       +\|\partial_R^4\tau_{R,h}\|_{C^0}\bigr)<C;
\]
\item the simple eigenprojector near \((u,s,b)=(0,0,0)\) is analytic in the
      twist and \(C^2\) in \(R\), with a uniform spectral gap.
\end{enumerate}
\end{theorem}

\begin{proof}
Finite horizon gives common lower and upper flight times and common invariant
cones.  Homogeneity strips give polynomial derivative losses which are
summable against the one-step expansion.  Parent labels keep a branch fixed
until its fold; the two one-sided traces are bounded by the stable-curve trace
theorem and their jump is exactly the fold contribution.  The growth lemma
and exponential return tail yield (ii) and (iii).  Compactness of the radius
interval gives a finite atlas.  Keller--Liverani perturbation on each
strong/weak pair, together with equality of physical pairings on overlaps,
gives (iv).
\end{proof}

\subsection{Arithmetic rank and returned UNI}

The interval-Newton certificate in
\texttt{A2\_ROUND11\_PERIODIC\_CERTIFICATE.json} contains five regular loops
with at least three collisions, clearance at least \(1.4\times10^{-3}\), and
incidence at least \(6.8\times10^{-2}\).  Their vectors
\((\kappa_1,\kappa_2,r)\) generate \(\mathbb Z^3\).  The same certificate
contains two returned branches with a temporal derivative gap at least
\(1.2\times10^{-2}\), while the common-suffix derivative is at most
\(2.5\times10^{-3}\).

\begin{lemma}[Uniform arithmetic obstruction]
\label{lem:r17-a2-arithmetic}
If
\[
 e^{i(u\cdot\kappa+s r+b\tau_R)}
 =e^{ic}{g\over g\circ F_R}
\]
on every regular return branch, then \((u,s)=(0,0)\) in
\(\mathbb T^3\) and \(b=0\).  The conclusion is uniform in \(R\).
\end{lemma}

\begin{proof}
Evaluation on the certified loops forces the three lattice characters to be
trivial because their integer vectors have determinant bounded away from
zero.  With \((u,s)=0\), evaluate on the two returned branches and
differentiate their phase difference along the certified unstable interval.
The temporal derivative gap is positive uniformly in \(R\), whereas a
coboundary has equal returned derivatives.  Thus \(b=0\).
\end{proof}

\begin{lemma}[Returned temporal nonintegrability]
\label{lem:r17-a2-uni}
There are two inverse return branches \(h_R^\pm\), a common unstable interval
\(J_R\), and \(c_{\rm UNI}>0\) such that
\[
 \inf_{R\in[0.45,0.47]}
 \inf_{x\in J_R}
 \left|\partial_x\bigl(\tau_R\circ h_R^+
                    -\tau_R\circ h_R^-\bigr)(x)\right|
 \ge c_{\rm UNI}.
\]
All branches and their first four derivatives remain inside the certified
incidence and clearance boxes.
\end{lemma}

\begin{proof}
The reflection equations and their Jacobian are evaluated by interval Newton
on each certified box.  The residual bound gives a unique orbit in the box;
the incidence and clearance bounds keep it regular for every radius.  The
interval derivative enclosure for the two branches is disjoint and its gap is
the stated certificate value.  Smooth continuation in \(R\) preserves the
same interval boxes.
\end{proof}

\subsection{Vertical smoothing of matrix coefficients}

For a regular branch word \(\mathbf h=(h_1,\ldots,h_n)\), use collision
arclength coordinates on each returned unstable curve.  The conditional
amplitude, including Jacobians and a central insertion, is denoted
\(a_{\mathbf h,R}\).  A block is \emph{vertical-regular} when it contains the
returned pair from Lemma~\ref{lem:r17-a2-uni}.  The growth lemma gives a
positive frequency of such blocks outside an exponentially small family.

\begin{lemma}[Raw branchwise integration by parts]
\label{lem:r17-a2-ibp}
For every \(M\ge2\) there are \(C_M,c_M,\kappa_M>0\) such that, uniformly in
\(R\), \(n\), central insertions in a fixed unit ball, and lattice twists,
\[
 |\Phi_{n,R}^{\ell}(u,s,b)|
 \le C_M\rho^n(1+|b|)^{-2}
      +C_Me^{c_Mn}(1+|b|)^{-M},
 \qquad |b|\ge1.
\]
The estimate is for the unsmoothed characteristic function itself.
\end{lemma}

\begin{proof}
Cut every branch word by a smooth partition supported strictly inside its
homogeneity components.  Boundary pieces are reexpressed by the two trace
components in Theorem~\ref{thm:r17-a2-bundle}; their total strong-to-weak norm
is summable.  On a vertical-regular block use the phase difference of the two
returned branches as one coordinate.  Lemma~\ref{lem:r17-a2-uni} supplies a
Jacobian bounded away from zero.  Integration by parts transfers one
\(b^{-1}\) factor to the conditional amplitude.  Repeating on disjoint blocks
transfers \(M\) factors.  Distortion and the branch-complexity sum bound all
derivatives by \(C_Me^{c_Mn}\).

Words with fewer than the required regular blocks have transfer weight at
most \(C\rho^n\).  They still contain two ordinary non-grazing flight
coordinates; two coarea integrations give the common factor
\((1+|b|)^{-2}\).  No frequency-independent exceptional term remains.
Summing the branch words by (iii) of
Theorem~\ref{thm:r17-a2-bundle} proves the estimate.
\end{proof}

\begin{theorem}[Exhaustive integrable Fourier bounds]
\label{thm:r17-a2-fourier}
There are \(\delta,c,C,A,\kappa>0\) such that the dual domain is covered by
four estimates:
\begin{enumerate}[label=(\alph*)]
\item for \(|(u,s,b)|\le\delta\),
\[
 \Phi_{n,R}^{\ell}
 =e^{nP_R(iu,is,ib)}
  \bigl(\ell_R(\Pi_{R,u,s,b}h_R)+O(\rho^n)\bigr),
\]
with a uniform cubic expansion of \(P_R\);
\item on compact lattice/roof minor arcs the modulus is at most \(Ce^{-cn}\);
\item if \(1\le|b|\le e^{\kappa n}\),
\[
 |\Phi_{n,R}^{\ell}(u,s,b)|
 \le C(1+|b|)^Ae^{-cn};
\]
\item if \(|b|>e^{\kappa n}\), Lemma~\ref{lem:r17-a2-ibp} holds with
      \(M\) chosen so that \((M-1)\kappa>c_M+2\kappa\).
\end{enumerate}
Consequently
\[
 \int_{\mathbb R}
 \sup_{(u,s)\in\mathbb T^3}
 |\Phi_{n,R}^{\ell}(u,s,b)|\,db<\infty
\]
and the integral outside a shrinking Gaussian neighborhood is
\(o(n^{-2})\), uniformly in \(R\).
\end{theorem}

\begin{proof}
Part (a) is analytic perturbation from
Theorem~\ref{thm:r17-a2-bundle}.  Part (b) follows from
Lemma~\ref{lem:r17-a2-arithmetic}, compactness, and upper semicontinuity of the
spectral radius.  The Dolgopyat pairing based on
Lemma~\ref{lem:r17-a2-uni} gives (c); choose \(\kappa\) so small that the
polynomial factor and the length of the interval are dominated by
\(e^{-cn/2}\).  For (d), integrate the two terms of
Lemma~\ref{lem:r17-a2-ibp}; the selected inequality on \(M\) makes the second
integral exponentially small.  This proves the global integrable majorant and
the final assertion.
\end{proof}

\subsection{Raw four-dimensional local limit theorem}

Let
\[
 Z_n=(S_n\kappa,S_nr,S_n\tau_R)
 \in\mathbb Z^2\times\mathbb Z\times\mathbb R.
\]
For a real tilt in a compact simple branch let \(m_R\) and \(\Sigma_R\) be
the mean and the full \(4\times4\) covariance.  Arithmetic obstruction and
UNI imply \(cI_4\le\Sigma_R\le CI_4\).

\begin{theorem}[Unsmoothed vector--return--roof density LLT]
\label{thm:r17-a2-llt}
The measure of \(Z_n\) has a density in its roof coordinate.  Uniformly for
\((k,m,t)-nm_R=O(\sqrt n)\),
\[
 p_{n,R}(k,m,t)
 ={\exp\{-\tfrac1{2n}(z-nm_R)^T\Sigma_R^{-1}(z-nm_R)\}
   \over (2\pi n)^2\sqrt{\det\Sigma_R}}
 \bigl(\nu_R^{\rm tilt}(\ell)+o(1)\bigr),
\]
where \(z=(k_1,k_2,m,t)\).  The error is uniform for central insertions in a
fixed H\"older unit ball and for \(R\in[0.45,0.47]\).
\end{theorem}

\begin{proof}
Use Fourier inversion on \(\mathbb T^3\times\mathbb R\).  On the Gaussian
region insert the cubic pressure expansion from
Theorem~\ref{thm:r17-a2-fourier}(a), scale by \(n^{-1/2}\), and apply
dominated convergence.  The covariance lower bound gives the Gaussian
dominator.  Parts (b)--(d) give a single integrable majorant on the complement
and an \(o(n^{-2})\) contribution.  The spectral projector surrounding the
central insertion converges to the stationary tilted pairing.  No smoothing
test is introduced at any stage.
\end{proof}

\begin{corollary}[Sharp clock conditioning]
\label{cor:r17-a2-conditioning}
Conditioning on exact homology and return count and on any roof interval gives
the conditional Gaussian interval mass obtained by integrating the density in
Theorem~\ref{thm:r17-a2-llt}.  In particular the local, central, and saturated
window regimes are the three corresponding Gaussian interval asymptotics,
and central-window path laws converge to the tilted two-sided law.
\end{corollary}

\begin{proof}
Apply the density theorem to numerator and denominator, integrate over the
roof interval, and divide.  The covariance lower bound keeps the denominator
uniformly positive in every stated regime.
\end{proof}

\begin{theorem}[Round-seventeen A2 closure]
\label{thm:r17-a2-main}
The Sinai family has one constructed parent-fold induced bundle, a certified
full lattice obstruction, returned UNI, an unsmoothed integrable Fourier
theorem based on matrix coefficients, and the raw four-dimensional density
LLT.  No trace-class or flat-trace identification is used.
\end{theorem}

\begin{proof}
Combine Theorems~\ref{thm:r17-a2-bundle},
\ref{thm:r17-a2-fourier}, and \ref{thm:r17-a2-llt}, together with
Lemmas~\ref{lem:r17-a2-arithmetic}, \ref{lem:r17-a2-uni}, and
\ref{lem:r17-a2-ibp}.
\end{proof}
